\section{Chapitre 5 : Trigonom\'etrie}

%% ===============================================================
\subsection{Exercice 30 : \'Equations trigonom\'etriques}

\textbf{A) QCM.}

\textbf{1.} On a $\cos x-\sqrt3\sin x=2\cos\left(x+\dfrac{\pi}{3}\right)$. L'\'equation $\cos\left(x+\frac{\pi}{3}\right)=0$ donne
$x=\dfrac{\pi}{6}+k\pi$. Dans $]-\pi;\pi]$ : $S=\left\{\dfrac{\pi}{6};-\dfrac{5\pi}{6}\right\}$.
Aucune proposition n'est exacte : \textbf{r\'eponse d)}. (La proposition c) $\{\frac{\pi}{6}\}$ est incompl\`ete.)

\textbf{2.} On retranche un multiple de $2\pi$ : $\dfrac{498\pi}{5}-50\times2\pi=\dfrac{498\pi-500\pi}{5}=-\dfrac{2\pi}{5}$, qui est bien dans $]-\pi;\pi]$.
\textbf{R\'eponse a)} $-\dfrac{2\pi}{5}$.

\textbf{3.} $\cos x=\dfrac{\sqrt6+\sqrt2}{4}=\cos\dfrac{\pi}{12}$, avec $x\in\left[0;\dfrac{\pi}{2}\right]$ : $x=\dfrac{\pi}{12}$.
Alors $\cos2x=\cos\dfrac{\pi}{6}=\dfrac{\sqrt3}{2}$ : \textbf{r\'eponse d)}.

\textbf{B) Vrai ou faux.}
\begin{enumerate}
  \item \textbf{Faux.} Une application surjective n'est pas forc\'ement injective : par exemple $x\mapsto x^2$ de $\R$ vers $\R_+$ est surjective mais pas injective.
  \item \textbf{Vrai.} $\cos x=-\frac12$ sur $]-\pi;\pi]$ donne $x=\pm\frac{2\pi}{3}$, soit $S=\left\{-\frac{2\pi}{3};\frac{2\pi}{3}\right\}$.
\end{enumerate}

%% ===============================================================
\subsection{Exercice 31 : \'Equations et in\'equations}

Sur $[0;2\pi]$ : $\sqrt3\cos2x-\sin2x=2\cos\left(2x+\dfrac{\pi}{6}\right)$. L'in\'equation (I) devient
\[
2\cos\left(2x+\tfrac{\pi}{6}\right)+1<0\iff\cos\left(2x+\tfrac{\pi}{6}\right)<-\frac12.
\]
Posons $\theta=2x+\frac{\pi}{6}$, avec $\theta\in\left[\frac{\pi}{6};\frac{25\pi}{6}\right]$. On a $\cos\theta<-\frac12$ sur
$\left]\frac{2\pi}{3}+2k\pi;\frac{4\pi}{3}+2k\pi\right[$.
\begin{itemize}
  \item $k=0$ : $\theta\in\left]\frac{2\pi}{3};\frac{4\pi}{3}\right[$, soit $2x\in\left]\frac{\pi}{2};\frac{7\pi}{6}\right[$, donc $x\in\left]\frac{\pi}{4};\frac{7\pi}{12}\right[$.
  \item $k=1$ : $\theta\in\left]\frac{8\pi}{3};\frac{10\pi}{3}\right[$, soit $2x\in\left]\frac{5\pi}{2};\frac{19\pi}{6}\right[$, donc $x\in\left]\frac{5\pi}{4};\frac{19\pi}{12}\right[$.
\end{itemize}
\textbf{$S=\left]\frac{\pi}{4};\frac{7\pi}{12}\right[\ \cup\ \left]\frac{5\pi}{4};\frac{19\pi}{12}\right[$} : \textbf{r\'eponse a)}
(\`a l'ouverture pr\`es des intervalles, qui sont ouverts car l'in\'egalit\'e est stricte).

%% ===============================================================
\subsection{Exercice 32 : \'Equations trigonom\'etriques}

\textbf{(E1)} $\cos x-\sin x=\sqrt2$. Comme $\cos x-\sin x=\sqrt2\cos\left(x+\frac{\pi}{4}\right)$ : $\cos\left(x+\frac{\pi}{4}\right)=1$,
donc $x=-\dfrac{\pi}{4}+2k\pi$, $k\in\mathbb Z$.

\textbf{(E2)} $\sqrt3\sin x+3\cos x=\sqrt3$. Comme $\sqrt3\sin x+3\cos x=2\sqrt3\sin\left(x+\frac{\pi}{3}\right)$ :
$\sin\left(x+\frac{\pi}{3}\right)=\frac12$, donc $x+\frac{\pi}{3}=\frac{\pi}{6}+2k\pi$ ou $\frac{5\pi}{6}+2k\pi$.
\[
S_{E_2}=\left\{-\frac{\pi}{6}+2k\pi\ ;\ \frac{\pi}{2}+2k\pi\ ;\ k\in\mathbb Z\right\}.
\]

\textbf{(E3)} $\sqrt3\sin x+\cos x=-1$. Comme $\sqrt3\sin x+\cos x=2\sin\left(x+\frac{\pi}{6}\right)$ :
$\sin\left(x+\frac{\pi}{6}\right)=-\frac12$, donc $x+\frac{\pi}{6}=-\frac{\pi}{6}+2k\pi$ ou $\frac{7\pi}{6}+2k\pi$.
\[
S_{E_3}=\left\{-\frac{\pi}{3}+2k\pi\ ;\ \pi+2k\pi\ ;\ k\in\mathbb Z\right\}.
\]
(V\'erification : pour $x=\pi$, $\sqrt3\cdot0+(-1)=-1$.)

%% ===============================================================
\section{Chapitre 6 : G\'eom\'etrie dans l'espace}

%% ===============================================================
\subsection{Exercice 33 : Cube et orthogonalit\'e}

On utilise le rep\`ere $\left(A;\vect{AB},\vect{AD},\vect{AE}\right)$ avec le cube de c\^ot\'e $a$ :
$A(0;0;0)$, $B(a;0;0)$, $C(a;a;0)$, $D(0;a;0)$, $E(0;0;a)$, $F(a;0;a)$, $G(a;a;a)$, $H(0;a;a)$.
(Figure : face $ABCD$ en bas, face $EFGH$ au-dessus, $E$ au-dessus de $A$.)
Les milieux : $I$ milieu de $[BF]$ : $I\left(a;0;\frac a2\right)$ ; $J$ milieu de $[AE]$ : $J\left(0;0;\frac a2\right)$.

\textbf{1.} Figure : voir le sch\'ema ci-dessus (cube, puis $I$ et $J$ sur les ar\`etes $[BF]$ et $[AE]$).

\textbf{2.} $\vect{ED}=(0;a;-a)$ et $\vect{BG}=(0;a;a)$. Leur produit scalaire vaut $a^2-a^2=0$.
\textbf{Donc $(ED)\perp(BG)$.}

\textbf{3.} $\vect{IJ}=\left(-a;0;0\right)$. Le plan $(ADE)$ est le plan $x=0$, de vecteur normal $(1;0;0)$ :
$\vect{IJ}$ est normal au plan. \textbf{Donc $(IJ)\perp(ADE)$.}

\textbf{4.} Plan $(EAC)$ : il contient $A$, $E$ et $C$, donc un vecteur normal est $\vect n_1=(1;-1;0)$
(orthogonal \`a $\vect{AE}=(0;0;a)$ et $\vect{AC}=(a;a;0)$).
Plan $(HDB)$ : $\vect{DB}=(a;-a;0)$ et $\vect{DH}=(0;0;a)$, donc un vecteur normal est $\vect n_2=(1;1;0)$.
$\vect n_1\cdot\vect n_2=1-1=0$. \textbf{Donc les plans $(EAC)$ et $(HDB)$ sont perpendiculaires.}

\textbf{5.} Plan $(DBF)$ : $\vect{DB}=(a;-a;0)$ et $\vect{DF}=(a;-a;a)$. Le vecteur $(1;1;0)$ est normal ($\vect{DB}\cdot(1;1;0)=0$ et
$\vect{DF}\cdot(1;1;0)=0$) ; le plan a pour \'equation $x+y=a$ (v\'erifi\'ee en $D$, $B$, $F$).
\[
d(A,(DBF))=\frac{|0+0-a|}{\sqrt{1^2+1^2}}=\frac{a}{\sqrt2}=\frac{a\sqrt2}{2}.
\]

%% ===============================================================
\section{Chapitre 7 : QCM et questions diverses}

%% ===============================================================
\subsection{Exercice 34 : QCM}

\textbf{A.1.} $f(x)=\dfrac{5x^3-3x-2}{x^4-x^3}$ (lecture de l'\'enonc\'e). Au voisinage de $+\infty$, $f(x)\sim\dfrac{5x^3}{x^4}=\dfrac5x\to0$.
La limite vaut $0$ : \textbf{r\'eponse d)} (aucune des trois autres).
\begin{rem} Si le num\'erateur est diff\'erent, la conclusion peut changer ; v\'erifier l'\'enonc\'e.\end{rem}

\textbf{A.2.} $g(x)=\dfrac{-2x^2-x+1}{x^2+1}$. Le d\'enominateur ne s'annule jamais : \textbf{pas d'asymptote verticale}.
En $\pm\infty$, $g(x)\to-2$ : la droite $y=-2$ est asymptote horizontale. \textbf{R\'eponse b).}

\textbf{B.1.} La tangente \`a $\mathcal C_f$ au point d'abscisse $a$ a pour \'equation $\boxed{y=f'(a)(x-a)+f(a)}$.

\textbf{B.2.} $\displaystyle\lim_{x\to0}\frac{\sin x}{x}=\mathbf{1}$.

\textbf{B.3.} Rotation de centre $O$ et d'angle $\frac{2\pi}{3}$ : $\cos\frac{2\pi}{3}=-\frac12$ et $\sin\frac{2\pi}{3}=\frac{\sqrt3}{2}$, donc
\[
\begin{cases}
x'=-\dfrac12x-\dfrac{\sqrt3}{2}y\\[6pt]
y'=\dfrac{\sqrt3}{2}x-\dfrac12y.
\end{cases}
\]

%% ===============================================================
\subsection{Exercice 35 : Vrai/Faux et compl\'ements}

\textbf{Partie A.}
\begin{enumerate}
  \item \textbf{Vrai.} Une fonction bijective est injective (et surjective).
  \item \textbf{Faux.} $\sin x\in[-1;1]$, donc $\sin x=\frac32$ n'a aucune solution.
\end{enumerate}

\textbf{Partie B.}
\begin{enumerate}
  \item Une fonction $f$ est une \textbf{application} si l'ensemble de d\'epart et l'ensemble de d\'efinition coïncident.
  \item Deux r\'eels de somme $S$ et de produit $P$ sont les solutions de $\boxed{x^2-Sx+P=0}$.
  \item $f:E\to F$ est \textbf{injective} si chaque \'el\'ement de $F$ admet au plus un ant\'ec\'edent.
  \item $x+y=\dfrac{2\pi}{3}-\dfrac{\pi}{4}=\dfrac{5\pi}{12}$, d\'ej\`a dans $]-\pi;\pi]$ : \textbf{mesure principale $\dfrac{5\pi}{12}$}.
  \item $h(M)=M'$ \'equivaut \`a $\boxed{\vect{\Omega M'}=-2\,\vect{\Omega M}}$.
  \item Si $G=\mathrm{bar}\{(A;\alpha);(B;\beta)\}$ avec $\alpha+\beta=1$, alors $\vect{BG}=\alpha\,\vect{BA}$.
  Or $\vect{BG}=-\dfrac23\vect{BA}$ donne $\alpha=-\dfrac23$ et $\beta=1-\alpha=\dfrac53$.
\end{enumerate}

%% ===============================================================
\subsection{Exercice 36 : QCM}

\textbf{I. Vrai ou faux.}
\begin{enumerate}
  \item \textbf{Vrai.} $f(x)=\sin\left(2x-\frac{\pi}{3}\right)$ a pour p\'eriode $\frac{2\pi}{2}=\pi$.
  \item \textbf{Faux.} $\cos x=\frac14$ et $\cos y=\frac12$ avec $x,y\in\left[0;\frac{\pi}{2}\right]$ : $y=\frac{\pi}{3}$, $\sin x=\frac{\sqrt{15}}{4}$.
  $\cos2y=-\frac12$, $\sin2y=\frac{\sqrt3}{2}$ et
  \[
  \cos(x-2y)=\cos x\cos2y+\sin x\sin2y=-\frac18+\frac{\sqrt{45}}{8}=\frac{3\sqrt5-1}{8}.
  \]
  L'\'enonc\'e donne $\frac{3\sqrt5+1}{8}$ : \textbf{l'affirmation est fausse}.
  \item \textbf{Faux.} $\vect{AB}=3\vect{CA}$ donne $\vect{AB}=-3\vect{AC}$ : $B$ est l'image de $C$ par l'homoth\'etie de centre $A$ et de rapport $-3$ (et non $3$).
  \item \textbf{Vrai.} $MA^2-MB^2=0\iff MA=MB$ : c'est la m\'ediatrice de $[AB]$.
\end{enumerate}

\textbf{II. Choix multiple} (\'enonc\'e non reproduit dans le texte).

%% ===============================================================
\subsection{Exercice 37 : QCM}

\textbf{I.1.} $3x^2-2|x|-1=0$. En posant $t=|x|\ge0$ : $3t^2-2t-1=0$, soit $(3t+1)(t-1)=0$, donc $t=1$ (car $t\ge0$).
Alors $x=\pm1$ : \textbf{deux solutions distinctes, r\'eponse b)}.

\textbf{I.2.} $\sqrt{x^2-2x}>1-x$ : domaine $x^2-2x\ge0$, soit $x\le0$ ou $x\ge2$.
\begin{itemize}
  \item Si $1-x<0$ (donc $x>1$), l'in\'egalit\'e est vraie d\`es que $x\ge2$.
  \item Si $1-x\ge0$ et $x\le0$ : on \'el\`eve au carr\'e, $x^2-2x>(1-x)^2=1-2x+x^2$, soit $0>1$ : impossible.
\end{itemize}
Donc $S=[2;+\infty[$ : \textbf{r\'eponse a)}.

\textbf{I.3.} L'\'enonc\'e \'ecrit $3\vect{BC}-4\vect{AC}=2\vect{BC}$, soit $\vect{BC}=4\vect{AC}$, ce qui s'\'ecrit $4A-B-3C=\vec 0$.
Cette relation n'est pas de la forme $\sum\alpha_i\vect{GX_i}=\vec 0$ avec les poids des options (les poids de l'\'enonc\'e ne somment pas \`a $0$) :
\begin{rem} L'\'enonc\'e est incoh\'erent (il manque probablement un point ou une lettre). Aucune des quatre r\'eponses ne se d\'eduit
directement de la relation \'ecrite ; merci de v\'erifier la formule originale.\end{rem}

%% ===============================================================
\subsection{Exercice 38 : Ajustement lin\'eaire}

\textbf{Consigne 1.} La production journali\`ere est $h(v)=-v+3300-\dfrac{8100}{v}$ pour $v\in[50;120]$.
\[
h'(v)=-1+\frac{8100}{v^2}=\frac{8100-v^2}{v^2}.
\]
$h'(v)=0\iff v=90$ (on prend $v>0$) ; $h'>0$ avant, $h'<0$ apr\`es : maximum en $v=90$.
$h(90)=-90+3300-90=3120$. \textbf{La production maximale est atteinte pour $v=90$, avec $3120$ exemplaires par jour.}

\medskip\noindent\textbf{Consigne 2.} Ajustement par moindres carr\'es de $y$ en fonction de $x$ ($x$ en milliers de FCFA) :
$\bar x=130$, $\bar y=504$, $\sum(x_i-\bar x)^2=16\,800$, $\sum(x_i-\bar x)(y_i-\bar y)=-100\,060$.
\[
a=\frac{-100\,060}{16\,800}\approx-5{,}956,\qquad b=\bar y-a\bar x\approx1\,278{,}27,\qquad \hat y=-5{,}956\,x+1\,278{,}27.
\]
\textbf{B\'en\'efice} : chaque exemplaire est vendu $1000x$ F et co\^ute $25\,000$ F ; les frais fixes sont $28\,000\,000$ F. Donc
\[
B(x)=\hat y\,(1000x-25\,000)-28\,000\,000.
\]
En d\'eveloppant : $B(x)=-5\,956\,x^2+1\,427\,170\,x-59\,956\,750$.
Le maximum est atteint pour $x=\dfrac{1\,427\,170}{2\times5\,956}\approx119{,}8$.
\textbf{Le prix de vente optimal est d'environ $119\,800$ F, pour un b\'en\'efice maximal d'environ $25{,}5$ millions de FCFA.}
\begin{hyp}le mod\`ele lin\'eaire $\hat y(x)$ est utilis\'e pour pr\'evoir le nombre d'exemplaires vendus.\end{hyp}

%% ===============================================================
\subsection{Exercice 39 : Probabilit\'es}

Une classe de $12$ filles et $9$ gar\c cons, soit $21$ \'el\`eves. Un comit\'e est un arrangement de $3$ \'el\`eves
(major, adjoint, secr\'etaire).

\textbf{1.} Nombre de comit\'es : $21\times20\times19=\mathbf{7980}$ : \textbf{r\'eponse b)}.

\textbf{2.} Comit\'es de m\^eme sexe : tous filles $12\times11\times10=1320$ ; tous gar\c cons $9\times8\times7=504$.
Total : $1320+504=\mathbf{1824}$ : \textbf{r\'eponse c)}.

\textbf{3.} $P(\text{que des filles})=\dfrac{1320}{7980}=\dfrac{22}{133}\approx0{,}165$ : \textbf{r\'eponse b)}.

%% ===============================================================
\subsection{Exercice 40 : Lancers de pi\`ece}

On lance deux fois une pi\`ece : $\Omega=\{(F;F);(F;P);(P;F);(P;P)\}$. Soit $A$ : ``obtenir au moins un pile'',
donc $A=\{(F;P);(P;F);(P;P)\}$.

\textbf{Appariement :}
\begin{itemize}
  \item ``Univers'' $\longleftrightarrow$ $\Omega=\{(F;F);(F;P);(P;F);(P;P)\}$ ;
  \item ``Obtenir uniquement les c\^ot\'es face'' $\longleftrightarrow$ l'\'ev\'enement $\{(F;F)\}$, qui est un \textbf{\'ev\'enement \'el\'ementaire} ;
  \item ``\'Ev\'enement incompatible avec $A$'' $\longleftrightarrow$ $\{(F;F)\}$, car $A\cap\{(F;F)\}=\varnothing$ ;
  \item ``\'Ev\'enement \'el\'ementaire'' $\longleftrightarrow$ $(P;P)$ ;
  \item ``\'Eventualit\'e'' $\longleftrightarrow$ $(F;P)$.
\end{itemize}
\begin{rem} L'\'enonc\'e propose une liste d'appariements dont l'un se recoupe : ``obtenir uniquement face'' est \`a la fois
\'ev\'enement \'el\'ementaire et incompatible avec $A$. La r\'eponse ci-dessus est la plus coh\'erente.\end{rem}

%% ===============================================================
\subsection{Exercice 41 : Cercle et repr\'esentation param\'etrique}

$x^2+y^2-12x+6y-4=0\iff(x-6)^2+(y+3)^2=4+36+9=49$.
\textbf{Centre $\Omega(6;-3)$, rayon $R=7$.}
\textbf{Repr\'esentation param\'etrique :} $\begin{cases}x=6+7\cos t\\ y=-3+7\sin t\end{cases}$, $t\in[0;2\pi[$.

%% ===============================================================
\subsection{Exercice 42 : Droites et angles}

Pour deux droites $(\Delta)$ et $(\Delta')$ s\'ecantes en $A$, la compos\'ee de deux sym\'etries d'axes concourants est une rotation de centre $A$ :
\[
S_{(\Delta')}\circ S_{(\Delta)}=\text{rotation de centre }A\text{ et d'angle }2\,\mathrm{Mes}(\vect u;\vect u')
=2\times\left(-\frac{\pi}{6}\right)=-\frac{\pi}{3}.
\]
\textbf{R\'eponse : la rotation de centre $A$ et d'angle $-\dfrac{\pi}{3}$.}

%% ===============================================================
\subsection{Exercice 43 : Compos\'ee de sym\'etries}

\textbf{La compos\'ee de deux sym\'etries d'axes parall\`eles est une \textbf{translation}} ; par contre la compos\'ee de deux
sym\'etries d'axes perpendiculaires est une \textbf{sym\'etrie centrale} (rotation d'angle $\pi$) de centre le point d'intersection des axes.

\section*{Conclusion}
Ces corrections couvrent les 43 exercices du recueil. Les points marqu\'es ``remarque'' ou ``hypoth\`ese'' signalent des
\'enonc\'es incomplets, incoh\'erents ou ambigus : il est conseill\'e de les v\'erifier aupr\`es du professeur avant de les utiliser.
